%!TEX root = ../csuthesis_research.tex

\keywordscn{持续学习；医学图像分类；语义分割；时间序列分类；多模态大语言模型}

\begin{abstractzh}

本调研报告围绕“基于闭式解算法的持续学习方法”展开系统分析。针对传统持续学习在增量阶段依赖梯度优化、容易产生灾难性遗忘，以及回放方法存在存储与隐私负担的问题，报告提出以解析学习为核心的技术路线，将增量更新转化为递归岭回归闭式解问题。该路线通过维护统计量并递推更新分类头，在不依赖历史原始样本的条件下实现知识保留与新知识吸收。

报告首先梳理持续学习研究现状与关键挑战，随后给出统一数学框架，包括问题定义、岭回归闭式解、递归更新公式及核心等价性证明。基于统一理论，报告进一步设计了面向医疗图像分类、时间序列分类、类增量语义分割和多模态专家路由的任务化方案，并给出对应框架图、模块设计与实验协议。实验部分从准确率、遗忘率、mIoU、后向迁移和增量耗时等维度提出评测标准，并规划了系统化的消融与鲁棒性分析流程。

在研究计划方面，报告制定了两个月（8周）推进方案，覆盖文献调研、算法实现、分任务实验、结果汇总与报告定稿等阶段，明确了每周任务与阶段性交付物。总体而言，本调研报告形成了“理论—方法—实验—计划”一体化的前置研究框架，可为后续课题深入实施提供完整依据。

\end{abstractzh}
